% Generated draft: requires mathematical and source review.
\documentclass[11pt]{article}
\usepackage[margin=25mm]{geometry}
\usepackage{fontspec}
\setmainfont{Latin Modern Roman}
\usepackage{amsmath,amssymb,mathtools,unicode-math}
\setmathfont{Latin Modern Math}
\usepackage{enumitem,needspace,xurl,hyperref}
\setlength{\parindent}{0pt}
\setlength{\parskip}{5pt}
\setlength{\emergencystretch}{3em}
\newcommand{\N}{\mathbb{N}}
\newcommand{\Z}{\mathbb{Z}}
\newcommand{\Q}{\mathbb{Q}}
\newcommand{\R}{\mathbb{R}}
\newcommand{\C}{\mathbb{C}}
\newcommand{\F}{\mathbb{F}}
\newcommand{\A}{\mathbb{A}}
\newcommand{\PP}{\mathbb{P}}
\newcommand{\E}{\mathbb{E}}
\newcommand{\eps}{\varepsilon}
\newcommand{\dd}{\,\mathrm{d}}
\begin{document}
\section*{UnsolvedMath ID: 800952}
\textbf{Source problem code:} catalogue entry
\subsection{Definitions and mathematical statement}
In the standard smooth, oriented formulation, let \(M\) be a closed connected surface of genus at least \(2\), let \(g\) be a Riemannian metric on \(M\), and let \(\sigma\) be its magnetic two-form. Define the Lorentz force \(Y:TM\to TM\) by
\[
 g_x(Y_xv,w)=\sigma_x(v,w).
\]
The magnetic flow \(\varphi^t\) consists of the trajectories satisfying
\[
 \nabla_{\dot\gamma}\dot\gamma=Y(\dot\gamma).
\]
Its kinetic energy and energy levels are
\[
 H(x,v)=\tfrac12 g_x(v,v),\qquad \Sigma_e=H^{-1}(e).
\]
For a positive energy \(e>0\), the level \(\Sigma_e\) is regular and compact. Assume that
\[
 h_{\mathrm{top}}\!\left(\varphi^1|_{\Sigma_e}\right)=0
\]
and that the magnetic geodesics on \(\Sigma_e\) have no conjugate points, understood in the magnetic sense on the selected energy level.

\textbf{Question (Paternain).} Must there exist constants \(\kappa,c\in\R\) such that
\[
 K_g(x)=\kappa\quad\text{and}\quad
 \sigma_x=c\,(\mathrm{dvol}_g)_x
 \qquad\text{for every }x\in M?
\]
Here \(K_g\) is the Gaussian curvature and \(\mathrm{dvol}_g\) is the area form for the chosen orientation. These two conclusions are precisely the source's stated meaning of local symmetry.

\textbf{Unresolved qualifications.} The source does not explicitly require positive or regular energy. The formulation above makes the usual positive-energy interpretation explicit; its restriction \(e>0\) requires review against the intended source convention. At \(e=0\), the kinetic-energy level consists of stationary trajectories, so the meaning and applicability of the no-conjugate-points hypothesis require separate specification. The supplied excerpt also does not define magnetic conjugate points or explicitly state its smoothness and orientability conventions; the intended magnetic variational convention must be confirmed.
\subsection{Short English statement}
Paternain asks whether a magnetic flow on a closed surface of genus at least \(2\), with zero topological entropy and no conjugate points on a chosen energy level, must have constant Gaussian curvature and a magnetic form equal to a constant multiple of the metric area form. The draft uses the positive-energy interpretation; the source leaves this restriction and the precise conjugate-point convention unstated.
\subsection{Sources}
\begin{itemize}
\item Local UnsolvedMath registry, permanent ID \(800952\), source problem code \texttt{catalogue entry}; supplied title: \emph{Open Problems and Questions about Geodesics --- Question 32 (Paternain)}. The registry labels the problem \texttt{open}; this is supplied status metadata, not an independently verified assessment.
\item Burns and Matveev, \emph{Open Problems and Questions about Geodesics} (2019, as identified in the supplied record), Question 32. Supplied source URL: \url{https://arxiv.org/abs/1308.5417}. The record attributes the question/context to Paternain and lists Burns and Matveev as proposers/collectors.
\item Supplied provenance records access on 2026-10-16 and extraction method \texttt{arxiv\_latex}. These are source-record provenance claims; no source visit or independent literature verification was performed for this draft.
\end{itemize}
\end{document}
